\subsection{Dynamic Typing in Assignment}

Python is dynamically typed in the sense that names do not have fixed declared types. The object has a type. The name is only a binding to an object.

This follows directly from the assignment model developed earlier:

\begin{verbatim}
name ---> object
\end{verbatim}

The name can later be rebound to another object. That other object may have a different type.

\subsubsection{Names Do Not Have Fixed Declared Types}

In C, a declaration such as:

\begin{verbatim}
int score;
\end{verbatim}

creates a variable whose type is fixed by the declaration. The storage location is meant to hold an integer value compatible with that declared type.

Python does not work like that.

\begin{verbatim}
score = 42

print(f"score: {score}")
print(f"type(score): {type(score)}")
\end{verbatim}

Possible output:

\begin{verbatim}
score: 42
type(score): <class 'int'>
\end{verbatim}

The name \texttt{score} is not an integer variable in the C sense. It is a name currently bound to an integer object.

The integer object has integer behavior:

\begin{verbatim}
score = 42

print("selected int methods:")
print(f"{'bit_length':<15} {'bit_length' in dir(score)}")
print(f"{'bit_count':<15} {'bit_count' in dir(score)}")
print(f"{'__add__':<15} {'__add__' in dir(score)}")

print()
print(f"score.bit_length(): {score.bit_length()}")
print(f"score + 8:          {score + 8}")
\end{verbatim}

Possible output:

\begin{verbatim}
selected int methods:
bit_length     True
bit_count      True
__add__        True

score.bit_length(): 6
score + 8:          50
\end{verbatim}

Those methods belong to the object reached through \texttt{score}. They do not belong to the spelling of the name.

If the same name is rebound to a string, the available behavior changes because the object changed:

\begin{verbatim}
score = "No soup for you."

print(f"score: {score}")
print(f"type(score): {type(score)}")

print()
print("selected str methods:")
print(f"{'upper':<15} {'upper' in dir(score)}")
print(f"{'replace':<15} {'replace' in dir(score)}")
print(f"{'split':<15} {'split' in dir(score)}")

print()
print(score.upper())
print(score.replace("soup", "debugger"))
\end{verbatim}

Possible output:

\begin{verbatim}
score: No soup for you.
type(score): <class 'str'>

selected str methods:
upper          True
replace        True
split          True

NO SOUP FOR YOU.
No debugger for you.
\end{verbatim}

The name is the same. The object is different. Therefore the type and behavior are different.

The practical rule is:

\begin{quote}
Do not ask what type a Python name has permanently. Ask what object the name currently references.
\end{quote}

\subsubsection{Rebinding the Same Name to Objects of Different Types}

The same Python name can be rebound to objects of different types during execution.

\begin{verbatim}
value = 42
print(f"{'step 1':<8} value={value!r:<30} type={type(value)}")

value = 3.14
print(f"{'step 2':<8} value={value!r:<30} type={type(value)}")

value = "D'oh!"
print(f"{'step 3':<8} value={value!r:<30} type={type(value)}")

value = True
print(f"{'step 4':<8} value={value!r:<30} type={type(value)}")

value = None
print(f"{'step 5':<8} value={value!r:<30} type={type(value)}")
\end{verbatim}

Possible output:

\begin{verbatim}
step 1   value=42                             type=<class 'int'>
step 2   value=3.14                           type=<class 'float'>
step 3   value="D'oh!"                        type=<class 'str'>
step 4   value=True                           type=<class 'bool'>
step 5   value=None                           type=<class 'NoneType'>
\end{verbatim}

This is legal because the name \texttt{value} is not declared as one fixed type. Each assignment changes the binding.

Conceptually:

\begin{verbatim}
step 1:
    value ---> int object 42

step 2:
    value ---> float object 3.14

step 3:
    value ---> str object "D'oh!"

step 4:
    value ---> bool object True

step 5:
    value ---> None singleton object
\end{verbatim}

This flexibility is useful in small scripts and interactive work, but it can also make code unclear if the same name changes meaning too often.

Poor style:

\begin{verbatim}
data = 42
data = "forty-two"
data = True
\end{verbatim}

Better:

\begin{verbatim}
answer_number = 42
answer_text = "forty-two"
answer_is_known = True
\end{verbatim}

The second version uses separate names for separate roles. Python does not require this, but human readers benefit from it.

A more realistic example:

\begin{verbatim}
raw_score = "42"
score = int(raw_score)
passed = score >= 40

print(f"raw_score: {raw_score!r}, type={type(raw_score)}")
print(f"score:     {score!r}, type={type(score)}")
print(f"passed:    {passed!r}, type={type(passed)}")
\end{verbatim}

Possible output:

\begin{verbatim}
raw_score: '42', type=<class 'str'>
score:     42, type=<class 'int'>
passed:    True, type=<class 'bool'>
\end{verbatim}

Here the names explain the transformation:

\begin{verbatim}
raw_score  -> original string text
score      -> converted integer
passed     -> boolean result of comparison
\end{verbatim}

This is clearer than repeatedly reusing one name for all stages.

\subsubsection{Runtime Type Inspection with \texttt{type()} and \texttt{isinstance()}}

Python allows type inspection at runtime.

The function \texttt{type()} returns the actual type object of a value:

\begin{verbatim}
items = [42, 3.14, "No soup for you.", True, None]

for item in items:
    print(f"{repr(item):<20} type={type(item)}")
\end{verbatim}

Possible output:

\begin{verbatim}
42                   type=<class 'int'>
3.14                 type=<class 'float'>
'No soup for you.'   type=<class 'str'>
True                 type=<class 'bool'>
None                 type=<class 'NoneType'>
\end{verbatim}

The function \texttt{isinstance()} asks whether an object belongs to a type or to a family of types.

\begin{verbatim}
value = 42

print(f"is int:     {isinstance(value, int)}")
print(f"is float:   {isinstance(value, float)}")
print(f"is str:     {isinstance(value, str)}")
print(f"is bool:    {isinstance(value, bool)}")
\end{verbatim}

Possible output:

\begin{verbatim}
is int:     True
is float:   False
is str:     False
is bool:    False
\end{verbatim}

A string example:

\begin{verbatim}
quote = "And now for something completely different."

print(f"type(quote): {type(quote)}")
print(f"is str:      {isinstance(quote, str)}")
print(f"is int:      {isinstance(quote, int)}")

print()
print("selected str methods:")
print(f"{'upper':<12} {'upper' in dir(quote)}")
print(f"{'lower':<12} {'lower' in dir(quote)}")
print(f"{'replace':<12} {'replace' in dir(quote)}")
\end{verbatim}

Possible output:

\begin{verbatim}
type(quote): <class 'str'>
is str:      True
is int:      False

selected str methods:
upper        True
lower        True
replace      True
\end{verbatim}

For simple exact checks, \texttt{type()} is direct:

\begin{verbatim}
value = 42

print(type(value) is int)
\end{verbatim}

Possible output:

\begin{verbatim}
True
\end{verbatim}

But \texttt{isinstance()} is usually better when inheritance matters.

The classic example is \texttt{bool}:

\begin{verbatim}
flag = True

print(f"type(flag): {type(flag)}")

print()
print(f"type(flag) is bool: {type(flag) is bool}")
print(f"type(flag) is int:  {type(flag) is int}")

print()
print(f"isinstance(flag, bool): {isinstance(flag, bool)}")
print(f"isinstance(flag, int):  {isinstance(flag, int)}")
\end{verbatim}

Possible output:

\begin{verbatim}
type(flag): <class 'bool'>

type(flag) is bool: True
type(flag) is int:  False

isinstance(flag, bool): True
isinstance(flag, int):  True
\end{verbatim}

This happens because \texttt{bool} is a subclass of \texttt{int}. The exact type is \texttt{bool}, but the object also passes an \texttt{int} family check.

\texttt{isinstance()} can also check against several possible types at once:

\begin{verbatim}
values = [42, 3.14, 2 + 5j, "D'oh!", True]

for value in values:
    is_number = isinstance(value, (int, float, complex))
    print(f"{repr(value):<10} numeric-like: {is_number}")
\end{verbatim}

Possible output:

\begin{verbatim}
42         numeric-like: True
3.14       numeric-like: True
(2+5j)     numeric-like: True
"D'oh!"    numeric-like: False
True       numeric-like: True
\end{verbatim}

The final line is true because \texttt{bool} is integer-compatible. Sometimes that is desirable. Sometimes it is not.

If booleans should be excluded, say so explicitly:

\begin{verbatim}
values = [42, 3.14, True, False]

for value in values:
    is_plain_number = (
        isinstance(value, (int, float)) and
        not isinstance(value, bool)
    )

    print(f"{repr(value):<8} plain number: {is_plain_number}")
\end{verbatim}

Possible output:

\begin{verbatim}
42       plain number: True
3.14     plain number: True
True     plain number: False
False    plain number: False
\end{verbatim}

The practical rule is:

\begin{quote}
Use \texttt{type()} when you need to inspect the exact runtime type. Use \texttt{isinstance()} when you need to ask whether an object belongs to an acceptable category of types.
\end{quote}

The structural summary for dynamic assignment is:

\begin{verbatim}
name = 42       -> name references an int object
name = "text"   -> same name now references a str object
name = None     -> same name now references the None singleton

type(name)      -> inspect the current object's exact type
isinstance(...) -> test whether the current object fits a type category
\end{verbatim}

Python assignment is dynamic because the name can move. Python objects remain typed because each object still carries its runtime type and behavior.